%% kapa: exam=1 round=26 year=2015 period=2 form=A subject=accounting subject_ko=회계학 seq=30 number=70
%% source: Q-Net 감정평가사 기출문제 (cst00309) artlSeq=5138001 file=extracted/1008/5138001/제26회_감정평가사 1차_기출/제26회_감정평가사 1차_기출/2교시_감평_A형.pdf
%% conversion: kordoc 4.24.51 -> tools/gen1.py (auto); review=proofread-sonnet; flags=none
\begin{question}{70}
(주)감평은 20×1년 중에 단기매매목적으로 공정가치 ￦10,000,000의 주식을 취득하고, 20×2년 7월 1일에 예외적인 상황이 발생하여 단기매매금융자산을 매도가능금융자산으로 재분류하였다. 동 주식의 공정가치는 다음과 같다.
\begin{framedtext}
\centering\begin{tabular}{ccc}
\uline{20×1년 말} & \uline{20×2년 7월 1일} & \uline{20×2년 말} \\
￦14,000,000 & ￦13,000,000 & ￦11,000,000
\end{tabular}\par
\end{framedtext}
동 주식과 관련하여 다음의 설명 중 옳지 않은 것은?
\begin{choices}
  \item 20×1년 순이익이 ￦4,000,000 증가한다.
  \item 20×2년 매도가능금융자산평가손실이 ￦2,000,000 인식된다.
  \item 20×2년 단기매매금융자산평가손실이 ￦1,000,000 인식된다.
  \item 20×2년 총자본이 ￦3,000,000 감소한다.
  \item 20×2년 기타포괄손익누계액이 ￦3,000,000 감소한다.
\end{choices}
\end{question}
\answer{5}
